% ID: SEL-D0035-OA-E04
% ESTADO: cotejada. Segunda pasada visual y borrador compilado revisados.
% Fuente: Selectividad/matematicas_aplicadas/titular_septiembre_EXAMEN 1-ANDALUCÍA.pdf
% SHA-256: 7a09a02641ea8de345b94f49b27515e58436e032c7401867f86dc523490a3ea5
% Página: 1; opción A; ejercicio original 4.

El tiempo de vida de una determinada especie de tortuga es una variable aleatoria que sigue una ley Normal de desviación típica 10 años. Se toma una muestra aleatoria simple de 10 tortugas y se obtienen los siguientes valores:
\[
46\quad38\quad59\quad29\quad34\quad32\quad38\quad21\quad44\quad34
\]
\begin{enumerate}[label=\alph*)]
\item \textbf{(1.5 puntos)} Determine un intervalo de confianza, al 95\%, para la vida media de dicha especie de tortugas.
\item \textbf{(1 punto)} Calcule el tamaño mínimo que debe tener una muestra para que el error de estimación de la vida media no sea superior a 5 años, con un nivel de confianza del 98\%.
\end{enumerate}
